\scanpage{115}
% archive foliation: 114

de sorte que $\unsure{\xi|\zeta}_1,\unsure{\xi|\zeta}_2,\unsure{\xi|\zeta}_3$ sont les solutions de l'équation
\[
(60)\qquad\underbrace{z^3+a_1z^2+a_2z+a_3}_{
\substack{\deleted{\unsure{(z-\unsure{\xi|\zeta}_j)(z-\unsure{\xi|\zeta}_h)}}\\\prod_i(z-\unsure{\xi|\zeta}_i)}}=0
\]
On aura, \deleted{donc} si
\[
\delta\in t_E=T_{E,0}\qquad\deleted{\unsure{\delta}}
\]
est une base de $t_E$, liée à $\wp$ mod $k.1$ par
\[
\varepsilon_2(\wp)\simeq\delta^{\otimes2}
\]
et si
\[
\wp'=\Theta_\delta\wp
\]
la relation
\[
\wp'^2=\wp^3+a_1\wp^2+a_2\wp+a_3.
\]

Déterminons l'autom involutif $\sigma_i$ de $\mathbb P^1_k$ qui échange
$\unsure{\xi|\zeta}_i$ et $\infty$, $\unsure{\xi|\zeta}_j$ et $\unsure{\xi|\zeta}_h$ --- qui \deleted{\unsure{est donc}}
\inserted{s'identifie à} l'autom déduit par passage au quotient de
\[
\deleted{x\longmapsto\unsure{x+}}\qquad x\longmapsto x+\varepsilon_i\qquad E\longrightarrow E.
\]
On trouve, \deleted{\unsure{en l'écrivant}}
\[
\sigma_i z=\frac{az+b}{cz+d}
\]
et écrivant $\sigma_i\infty=\unsure{\xi|\zeta}_i$, $\sigma_i\unsure{\xi|\zeta}_i=\infty$
\[
\deleted{\unsure{\sigma_i z=}}\qquad\frac ac=\unsure{\xi|\zeta}_i,\quad-\frac dc=\unsure{\xi|\zeta}_i\qquad\unsure{;} 
\]
soit, faisant $c=1$
\[
\sigma_i z=\frac{\unsure{\xi|\zeta}_i z+b}{\deleted{\unsure{\unsure{\xi|\zeta}_i}}z\unsure{-}\unsure{\xi|\zeta}_i}
\]
et écrivant $\sigma_i\unsure{\xi|\zeta}_j=\unsure{\xi|\zeta}_h$ on trouve
\[
\frac{\unsure{\xi|\zeta}_i\unsure{\xi|\zeta}_j+b}{\unsure{\xi|\zeta}_j-\unsure{\xi|\zeta}_i}=\unsure{\xi|\zeta}_h\quad\text{i.e.}\quad
b=\unsure{\xi|\zeta}_j\unsure{\xi|\zeta}_h-\unsure{\xi|\zeta}_i(\unsure{\xi|\zeta}_j+\unsure{\xi|\zeta}_h)
\]
